\scanpage{8}
\deleted{non divisible par $p$} alors \deleted{$Df$} $Df$ a en $x$ un pôle d'ordre au plus $n+1$, et égal à $n+1$ ssi $p\nmid n$.

\deleted{\uline{Corollaire 1} Si $k(s)$ est de car nulle, alors $\underline{\mathcal L}\otimes\underline{\mathcal L}_n\simeq\underline{\mathcal L}_{n+1}$ pour tt $n\geq2$, donc \deleted{$\underline{\mathcal L}_3\otimes\mathcal L\otimes\underline{\mathcal L}_2,\mathcal L_4$} \inserted{$\mathcal L_4\simeq\mathcal L_2\otimes\mathcal L^2$} $\mathcal L_n\simeq\mathcal L^{n-2}\otimes\mathcal L_2$, \deleted{donc} pour $n\geq2$, \deleted{\unsure{donc on} $\mathcal L_2\unsure{\simeq}$}}

\deleted{\uline{Corollaire} 2 Si $k(s)$ est de car $p\ne2$, alors \deleted{$\mathcal L_3\unsure{\otimes}\simeq\mathcal L_2\otimes\mathcal L$}, $\underline{\mathcal L}_4\simeq\underline{\mathcal L}_3\otimes\mathcal L$, $\mathcal L_6\simeq\underline{\mathcal L}_5\otimes\underline{\mathcal L}$ etc. $[\underline{\mathcal L}_{2n}=\underline{\mathcal L}_{2n-1}\otimes\mathcal L\ \text{si }n\geq2]$, enfin
\[
\mathcal L_2^{\deleted{\unsure{2}}}\simeq\mathcal L^{\otimes2}
\]}

\uline{Corollaire} Si $S$ est de car $p>0$ (i.e. $p\underline O_S=0$) et si $n$ est un multiple de la caractéristique, alors $\mathcal L\otimes\underline{\mathcal L}_n\longrightarrow\underline{\mathcal L}_{n+1}$ est \uline{nul}.

\deleted{\unsure{Propo|Prop|Prop.}} Considérons la suite exacte canonique
\[
\unsure{0|\cdot}\longrightarrow\underline L^{\otimes n}\longrightarrow\underline L^{\otimes n+1}\longrightarrow
(\underline{\mathfrak g}_{C/S}^{\otimes(n+1)})\otimes_{\underline O_{\unsure{C|S}}}\underline O_{S'}\longrightarrow0
\]
on en conclut la suite exacte
\[
0\longrightarrow\mathcal E_n\longrightarrow\mathcal E_{n+1}\longrightarrow\mathcal L^{\otimes(n+1)}\longrightarrow
\underset{\substack{\Vert\\0\ \text{si }n\geq1}}{\deleted{\unsure{R}}R^{1}(f,\underline L^{\otimes n})}
\]
d'où
\[
\boxed{\mathcal L_{n+1}\simeq\mathcal L^{\otimes(n+1)}\qquad n\geq2}
\]
% slanted
\margin{intervertir avec la prop. 2}
et la multiplication dont il est question dans prop. 2 n'est autre que la multiplication triviale pour cette forme des $\mathcal L_n$, $n\geq2$
% archive sheet number: 7
